\documentclass{beamer}
\usepackage{ctex}
\usetheme{SimplePlus}
\setbeamertemplate{footline}{}

\usepackage{tikz}
\usetikzlibrary{arrows.meta}
\usepackage{tikz-cd}

\usepackage[
  hybrid,                            % 允许在 Markdown 中混用 LaTeX 命令
  texMathDollars                     % 启用 $...$ 和 $$...$$ 数学公式
]{markdown}

\usepackage{unicode-math}
%\setmathfont{STIX Two Math}
%\setmathfont{TeX Gyre Termes Math}
%\setmathfont{Libertinus Math}
%\setmathfont{Fira Math}

%%%%%%%% FONT TIMESAVES %%%%%%%
\newcommand{\bT}{\mathbb{T}}
\newcommand{\bN}{\mathbb{N}}
\newcommand{\bI}{\mathbb{I}}

\newcommand{\CC}{\mathcal{C}}

%%%%%%%% CATEGORY THEORY %%%%%%%
\newcommand{\id}{\mathrm{id}}
\DeclareMathOperator{\Hom}{Hom}

\newcommand{\cat}[1]{\mathrm{#1}}
\newcommand{\Set}{\cat{Set}}

\DeclareFontFamily{U}{dmjhira}{}
\DeclareFontShape{U}{dmjhira}{m}{n}{ <-> dmjhira }{}

\DeclareRobustCommand{\yo}{\text{\usefont{U}{dmjhira}{m}{n}\symbol{"48}}}

\DeclareMathOperator*{\colim}{colim}

\newcommand{\inl}{\operatorname{inl}}
\newcommand{\inr}{\operatorname{inr}}

%%%%%%%%% GEOMTERY %%%%%%%%%%%%%%
\newcommand{\Spec}{\operatorname{Spec}}
\newcommand{\Zar}{\mathrm{Zar}}

%%%%%%%%% TYPE THEORY %%%%%%%%%%%
\newcommand{\OO}{\bigcirc}
\newcommand{\UU}{\mathcal{U}}
\newcommand{\isEquiv}{\bibopenparen{isEquiv}}
\newcommand{\hlevel}{\mathrm{hlevel}}
\newcommand{\isContr}{\operatorname{isContr}}
\newcommand{\isProp}{\operatorname{isProp}}
\newcommand{\Prop}{\mathrm{Prop}}
\newcommand{\refl}{\mathrm{refl}}
\newcommand{\isModal}{\operatorname{isModal}}
\newcommand{\ap}{\operatorname{ap}}
\newcommand{\fib}{\operatorname{fib}}
\newcommand{\im}{\operatorname{im}}
\newcommand{\tot}{\mathrm{tot}}
\newcommand{\pt}{\mathrm{pt}}

\title{cardinal(基数)}

\date{}

\begin{document}

\begin{frame}
    % Print the title page as the first slide
    \titlepage
\end{frame}

\begin{frame}[fragile]{严格全序 (Strict Total Order)}
\begin{markdown}
设 $A$ 为一个集合，$<$ 是 $A$ 上的二元关系。若 $<$ 满足以下三个性质，则称 $<$ 为 $A$ 上的**严格全序**（或严格线序）：

1. **非自反性 (Irreflexivity)**：对任意 $a \in A$，均有 $a \nless a$（即 $\neg(a < a)$）。
2. **传递性 (Transitivity)**：对任意 $a, b, c \in A$，若 $a < b$ 且 $b < c$，则 $a < c$。
3. **三歧性 (Trichotomy)**：对任意 $a, b \in A$，$a < b$、$a = b$、$b < a$ **恰有其一**成立。
\end{markdown}
\end{frame}

\begin{frame}[fragile]{严格良序 (Strict Well-ordering)}
\begin{markdown}
设 $A$ 为一个集合，$<$ 是 $A$ 上的二元关系。称 $(A, <)$ 为一个**严格良序集**（或 $<$ 是 $A$ 上的**严格良序**），当且仅当满足以下两条：

1. **关系 $<$ 是 $A$ 上的严格全序**。
2. **良基性 / 最小元原理 (Well-foundedness)**：$A$ 的**每一个非空子集**都有关于 $<$ 的最小元。
* 精确逻辑表达为：

$$\forall S \subseteq A \; (S \neq \emptyset \implies \exists m \in S \; \forall x \in S \setminus \{m\}, \; m < x)$$
\end{markdown}
\end{frame}

\begin{frame}{传递集 (Transitive Set)}
集合 $x$ 称为传递集，如果对任意 $y \in x$，都有 $y \subseteq x$（即 $x$ 的元素也是 $x$ 的子集）。
\end{frame}

\begin{frame}[fragile]{序数 (Ordinal)}
\begin{markdown}
按照冯·诺伊曼（von Neumann）的定义，一个集合 $\alpha$ 称为**序数**，当且仅当 $\alpha$ 是一个传递集，且在其上的属于关系 $\in$ 是严格良序的。
* 所有序数的类记作 $\mathbf{ON}$。
* 序数的大小关系定义为：$\alpha < \beta \iff \alpha \in \beta$。
\end{markdown}
\end{frame}

\begin{frame}[fragile,shrink]{极限序数（Limit Ordinal）}
\begin{markdown}
在集合论中，**极限序数（Limit Ordinal）** 是序数（Ordinal）分类中的一种特殊序数。直观地说，它是一个**无法通过对其紧前趋加 1 得到，而是由所有比它小的序数“取极限（并集）”拼凑而成**的无穷序数。

设 $\lambda$ 为一个序数（按照冯·诺伊曼定义，即一个传递集且在属于关系 $\in$ 下为严格良序集）。

$\lambda$ 称为**极限序数**，当且仅当它满足以下**任一等价条件**：

1. 标准定义（非零且非后继）
$$\lambda \neq 0 \quad \text{且} \quad \forall \alpha \, (\lambda \neq \alpha + 1)$$


即：$\lambda$ 既不是零序数 $0$（空集 $\emptyset$），也不能表示为任何序数 $\alpha$ 的后继序数 $\alpha + 1 = \alpha \cup \{\alpha\}$。

2. 局部最大元条件（无最大元）
$$\lambda \neq 0 \quad \text{且} \quad \forall \alpha < \lambda, \; \exists \beta < \lambda \text{ 使得 } \alpha < \beta$$


即：$\lambda$ 是非零的，且在比 $\lambda$ 小的所有序数组成的集合中，**没有最大元素**。

3. 并集/上确界条件（等于小序数的并）
$$\lambda \neq 0 \quad \text{且} \quad \lambda = \bigcup \lambda = \sup_{\alpha < \lambda} \alpha$$


即：$\lambda$ 恰好等于所有严格小于它的序数的并集（上确界）。
\end{markdown}
\end{frame}

\begin{frame}[fragile]{共尾数（Cofinality）}
\begin{markdown}
共尾数是衡量一个序数（或基数）“可以被多快逼近”的概念，是研究正则基数、奇异基数的核心工具。

1. 定义（最常用的形式）

设 $\alpha$ 是一个极限序数（或任意无限序数）。  
$\alpha$ 的**共尾数** $\operatorname{cf}(\alpha)$ 定义为：

> 最小的序数 $\beta$，使得存在一个函数  
> $f:\beta\to\alpha$，  
> 满足：$f$ 的值域在 $\alpha$ 中是**共尾的**（cofinal），  
> 即对任意 $\xi<\alpha$，都存在 $\eta<\beta$ 使得 $f(\eta)\ge\xi$。

换句话说，可以用长度为 $\beta$ 的递增序列“爬到” $\alpha$ 的顶端。

2. 等价刻画（更方便使用）

以下条件等价：
1. $\operatorname{cf}(\alpha)=\beta$
2. $\beta$ 是最小的序数，使得存在一个**严格递增**的函数 $f:\beta\to\alpha$，其值域在 $\alpha$ 中无上界。
3. $\beta$ 是最小的序数，使得 $\alpha$ 可以写成 $\beta$ 个比 $\alpha$ 小的序数的并：
   $$
   \alpha=\bigcup_{\xi<\beta}\alpha_\xi,\quad\text{其中每个 }\alpha_\xi<\alpha.
   $$
\end{markdown}
\end{frame}

\begin{frame}[fragile]{等势 (Equinumerosity)}
\begin{markdown}
* 设 $A, B$ 为两个集合。若存在一个从 $A$ 到 $B$ 的**双射**（一一对应） $f: A \to B$，则称 $A$ 与 $B$ **等势**，记作 $A \approx B$。
\end{markdown}
\end{frame}

\begin{frame}[fragile]{基数 (Cardinal Number)}
\begin{markdown}
在选择公理 (AC) 成立的条件下，任何集合都可以被良序化，因此可以将基数定义为一种特殊的序数。

\smallskip
序数 $\kappa$ 称为**基数 (Cardinal Number)**，当且仅当它**不能与任何比它小的序数等势**：
$$\kappa \text{ 是基数 } \iff \forall \alpha < \kappa, \; \alpha \not\approx \kappa$$

* **有限基数**：自然数 $0, 1, 2, \dots$
* **无限基数**：通常按序数下标表示为阿列夫数（$\aleph_0, \aleph_1, \aleph_2, \dots, \aleph_\omega, \dots$）。
* 每个集合 $A$ 的势 $\vert{}A\vert{}$ 即为与 $A$ 等势的唯一那个基数序数。
\end{markdown}
\end{frame}

\begin{frame}[fragile]{阿列夫序列 (Aleph Sequence)}
\begin{markdown}
在选择公理 (AC) 下，所有的无限基数都可以按照大小关系构成一个在序数（Ordinal）上索引的严密序列，称为**阿列夫序列**：
$$\aleph_0 < \aleph_1 < \aleph_2 < \dots < \aleph_\omega < \aleph_{\omega+1} < \dots$$

其中，每个阿列夫数 $\aleph_\alpha$ 的下标 $\alpha$ 都是一个序数。基数的类型（后继还是极限），**完全取决于其下标序数 $\alpha$ 的类型**。
\end{markdown}
\end{frame}

\begin{frame}[fragile]{ 后继基数 (Successor Cardinal)}
\begin{markdown}
1. 精确定义

一个无限基数 $\kappa$ 称为**后继基数**，当且仅当存在一个基数 $\mu < \kappa$，使得 $\kappa$ 是**严格大于 $\mu$ 的最小基数**。

* 通常记作：$\kappa = \mu^+$（读作 $\mu$ 的后继基数）。
* 若用阿列夫序列表示：若下标序数是后继序数（即 $\alpha = \beta + 1$），则 $\aleph_\alpha = \aleph_{\beta+1}$ 是后继基数。

2. 典型例子

* **$\aleph_1$**：$\aleph_0$（可数无穷）的直接下一个基数，即 $\aleph_1 = \aleph_0^+$（首个不可数基数）。
* **$\aleph_2, \aleph_3, \dots, \aleph_{100}$**：自然数下标的基数全部是后继基数。
* **$\aleph_{\omega+1}$**：极限序数 $\omega$ 的下一个后继序数对应的基数。

3. 核心性质

* **绝对正则性**：在 ZFC 集合论中，**所有后继基数都是正则基数 (Regular Cardinal)**。
也就是说，你不可能用少于 $\mu^+$ 个、且每个集合大小都严格小于 $\mu^+$ 的小集合通过并集拼凑出 $\mu^+$。
\end{markdown}
\end{frame}

\begin{frame}[fragile,shrink]{极限基数 (Limit Cardinal)}
\begin{markdown}
极限基数进一步分为**普通极限基数**（或简称极限基数）与**强极限基数**。

1. (普通) 极限基数 (Limit Cardinal)

一个无限基数 $\kappa$ 称为**极限基数**，当且仅当它满足以下两个条件：

1. $\kappa$ 不是有限基数，也不是 $\aleph_0$（即 $\kappa > \aleph_0$）；
2. $\kappa$ **不是任何基数的后继基数**（即对任意 $\mu < \kappa$，都有 $\mu^+ < \kappa$）。

若用阿列夫序列表示：若下标序数 $\alpha$ 是一个**极限序数**（如 $\omega, \omega\cdot 2, \omega_1$），则 $\aleph_\alpha$ 就是极限基数。

精确表达（上确界条件）

极限基数也可以表示为所有严格小于它的基数的上确界（并集）：
$$\kappa = \sup \{ \mu \mid \mu < \kappa \text{ 且 } \mu \text{ 是基数} \} = \bigcup \{ \mu \mid \mu < \kappa \text{ 且 } \mu \text{ 是基数} \}$$

典型例子

* **$\aleph_\omega$**：阿列夫序列中首个极限基数，它是 $\aleph_0, \aleph_1, \aleph_2, \dots$ 的上确界。
* **$\aleph_{\omega+\omega}$**（或 $\aleph_{\omega\cdot 2}$）、**$\aleph_{\omega_1}$**（下标为首个不可数序数）。

**核心性质**

* **多为奇异基数**：许多常见的极限基数（如 $\aleph_\omega$）都是**奇异基数 (Singular Cardinal)**，因为它们可以被严格小于自身的序列拼凑而成（例如 $\text{cf}(\aleph_\omega) = \aleph_0 < \aleph_\omega$）。（若兼具正则性则升级为弱不可达基数）
\end{markdown}
\end{frame}

\begin{frame}[fragile,shrink]
\begin{markdown}
2. 强极限基数 (Strong Limit Cardinal)

在涉及幂集运算时，还有一个更强烈的极限概念：

无限基数 $\kappa$ 称为**强极限基数**，当且仅当对任意基数 $\mu < \kappa$，其幂集基数均满足：
$$2^\mu < \kappa$$

与普通极限基数的区别

* **强极限条件更严苛**：不仅要求“后继运算” $\mu^+$ 跨不出 $\kappa$，还要求极具膨胀性的“幂集运算” $2^\mu$ 也**无法跨出 $\kappa$**。
* **关系**：因为在 ZFC 中总是成立 $\mu^+ \le 2^\mu$，所以**任意强极限基数必然是极限基数**；反之不一定（除非连续统假设 GCH 成立）。
* **典型代表**：
* **$\aleph_0$**：通常也被视为强极限（因为对有限数 $n$，有 $2^n < \aleph_0$）。
* **不可达基数**：强不可达基数的核心条件之一就是要求其为一个“强极限基数”。
\end{markdown}
\end{frame}

\begin{frame}[fragile]{正则基数 (Regular Cardinal)}
\begin{markdown}
正则基数用来刻画那些“无法被更小规模的运算从下方充填达到”的基数。

一个无限基数 $\kappa$ 称为**正则基数 (Regular Cardinal)**，当且仅当它的共尾性等于它自身：
$$\text{cf}(\kappa) = \kappa$$

若无限基数满足 $\text{cf}(\kappa) < \kappa$，则称其为**奇异基数 (Singular Cardinal)**。

\medskip
**直观等价性质**

$\kappa$ 是正则基数等价于以下判定：

> 不能将 $\kappa$ 表示为少于 $\kappa$ 个集合的并集，且其中每个集合的基数都严格小于 $\kappa$。
> 即：若 $\lambda < \kappa$，且对每个 $\alpha < \lambda$ 有 $\vert{}A_\alpha\vert{} < \kappa$，则 $\left\vert{} \bigcup_{\alpha < \lambda} A_\alpha \right\vert{} < \kappa$。
\end{markdown}
\end{frame}

\begin{frame}[fragile]{不可达基数 (Inaccessible Cardinal)}
\begin{markdown}
“不可达”意味着该基数极其庞大，以至于通过 ZFC 中标准的构造性操作（如取幂集、取小于该规模的并集）都无法从比它小的基数到达它。

数学界通常将不可达基数分为**弱不可达**和**强不可达**。如果不加特殊说明，学术界所说的“不可达基数”默认指**强不可达基数**。

* 弱不可达基数 (Weakly Inaccessible Cardinal)

无限基数 $\kappa$ 称为**弱不可达基数**，当且仅当它满足以下三个条件：

1. $\kappa > \aleph_0$（不可数）；
2. $\kappa$ 是**正则基数**；
3. $\kappa$ 是**极限基数**（即不是后继基数）。

* 强不可达基数 (Strongly Inaccessible Cardinal / 不可达基数)

无限基数 $\kappa$ 称为**强不可达基数**（通常直接称为**不可达基数**），当且仅当它满足以下三个条件：

1. $\kappa > \aleph_0$（不可数）；
2. $\kappa$ 是**正则基数**；
3. $\kappa$ 是**强极限基数**（即对任意 $\mu < \kappa$，都有 $2^\mu < \kappa$）。

\end{markdown}
\end{frame}


% 参考文献帧，长列表自动分页
\begin{frame}[allowframebreaks]{参考文献}
\begin{thebibliography}{99}

\bibitem{gemini}
\url{https://share.gemini.google/dYwMrvRcaywO}

\end{thebibliography}
\end{frame}

\end{document}
